\section{Current Status \& Next Steps}
\label{sec:conclusion}

\subsection{Where Things Stand}

The prototype is fully working, and we have completed an initial round of automated experiments (Section~\ref{sec:preliminary_results}). Here's a rundown of what's done:

\begin{itemize}
    \item Complete 7-stage quest with 16 static checkpoints + dynamic checkpoint generation
    \item 6 NPC Q-learning agents with pre-training, cold-start fallback, and live learning
    \item \textbf{Decomposed reward model:} penalty + individual + $\lambda \cdot$ community, with non-linear community reward (sigmoid reputation, sub-linear health, sigmoid mood)
    \item \textbf{Adaptive NPC personalities:} cooperation\_tendency, risk\_aversion, social\_sensitivity, shock\_resilience --- all evolving per turn based on reward feedback
    \item \textbf{Dynamic $\lambda$ feedback loop:} cooperation drives $\lambda$, which drives community reward weighting, creating virtuous/vicious cycles
    \item \textbf{Role-specific action masking:} soft Q-value adjustment per archetype, with configurable bonus/penalty
    \item \textbf{System shock engine:} 5 shock types with lifecycle management (activation $\to$ decay $\to$ expiry), stat drain, trust modifiers, reward scaling
    \item Provider-based LLM integration (4 use cases, 9-stage guardrails, template fallbacks) --- supports Ollama, llama.cpp, OpenAI-compatible endpoints
    \item \textbf{Research analytics dashboard:} Chart.js-powered cooperation curves, reward decomposition, social welfare, action distributions, shock timeline, experiment bundle export
    \item Real-time MDP graph in the browser via Cytoscape.js
    \item Save/load with auto-save, backups, corruption recovery, shock state serialization
    \item 28 universal actions (including \texttt{present\_item})
    \item Combat, per-NPC reputation, gossip propagation, stamina, difficulty presets, random events
    \item Dual input (button palette + free-text NLP parsing with optional LLM refinement)
    \item Research tools: playthrough logger (with RL telemetry), replay, playtest bot
    \item \textbf{91 tests} across 5 suites (core gameplay, shock engine, adaptation \& masking, analytics, integration)
    \item \texttt{uv} package management with \texttt{pyproject.toml}
\end{itemize}

\subsection{Key Design Decisions (Why We Did It This Way)}

A few of the non-obvious choices worth remembering:

\begin{itemize}
    \item \textbf{MDP owns the state, LLM just generates text.} The LLM never modifies game state directly --- everything goes through validation and the game engine. This prevents hallucination-induced state corruption.

    \item \textbf{All 28 actions always available.} We never grey out or block actions. If an action doesn't make sense in context, it produces a narrated failure instead. This keeps the action space observable for research.

    \item \textbf{LLM is always optional.} Every code path that calls the LLM has a template fallback. You can run the entire game with no model loaded. This makes ablation testing clean.

    \item \textbf{Tabular Q-learning on purpose.} We chose tabular over deep RL because you can literally open the Q-table and read what the NPC has learned. Full interpretability. With 225 states $\times$ 28 actions = 6,300 Q-values per NPC, the table is small enough to inspect directly.

    \item \textbf{Everything seeded.} All randomness comes from \texttt{MASTER\_SEED}. Same seed = same game. NPC pre-training seeds are derived deterministically from registry index, not \texttt{hash()}, ensuring cross-run reproducibility. Verified by explicit seed reproducibility tests.

    \item \textbf{Decomposed reward for ablation.} The penalty/individual/community split was chosen specifically so each component can be independently measured and turned off. The additive structure with dynamic $\lambda$ means we can study how the balance between selfish and prosocial utility shifts over time.

    \item \textbf{Non-linear community reward.} Sigmoid and sub-linear shaping prevent reward saturation at high welfare levels while amplifying sensitivity during crises --- the exact regime where cooperation should matter most.

    \item \textbf{Dynamic $\lambda$ over static.} A static $\lambda$ forces a fixed cooperation level. Dynamic $\lambda$ creates emergent feedback loops: positive community outcomes $\to$ cooperation rises $\to$ $\lambda$ increases $\to$ more prosocial behavior. This is the mechanism that produces the dynamics our research hypothesis predicts.

    \item \textbf{Shock engine for non-stationarity.} Rather than just relying on multi-agent interaction for non-stationarity, we added explicit environmental shocks that create controlled perturbations. This lets us study adaptation and recovery in a repeatable way.
\end{itemize}

\subsection{Known Limitations}

Things we're aware of but haven't fully addressed:

\begin{itemize}
    \item \textbf{Small scale:} 5 locations, 6 NPCs, 1 linear quest. We don't know if the architecture scales to bigger worlds or branching quest graphs.
    \item \textbf{Multi-agent convergence:} 6 Q-learners in a shared non-stationary environment --- no convergence guarantees \cite{bowling2002multiagent}. We see interesting emergent behavior but can't prove it's stable.
    \item \textbf{LLM quality:} Local models produce noticeably weaker text than cloud models. The guardrails have to work harder. Provider-based adapter helps (can point to better models), but quality varies.
    \item \textbf{No player modeling:} We don't track or predict the player's strategy. Adding inverse RL or belief tracking could help with adaptation.
    \item \textbf{Single session only:} One player at a time, no concurrent sessions.
    \item \textbf{Adaptation rate tuning:} The adaptation rate has been tuned from 0.04 to 0.018, with systematic role-specific multipliers (e.g., guard risk aversion $\times 1.5$, elder social sensitivity $\times 1.4$). Sensitivity to the base rate remains an open question for longer sessions.
    \item \textbf{Policy entropy:} Per-action counts are now tracked, providing full action distribution entropy ($H = -\sum_a p_i(a) \log p_i(a)$) rather than the previous 2-bin approximation.
\end{itemize}

\subsection{Key Experimental Findings}

The initial automated experiments (500-turn bot playthroughs, seed=42, quest auto-restart) have produced several quantitative findings that support the system's core design hypotheses:

\begin{itemize}
    \item \textbf{Cooperation emerges through RL:} The cooperation index $\cK$ rises from the baseline $0.500$ to $0.851$ over 500 turns with Q-learning (C1), while schedule-only NPCs (C3) remain at $0.500$ throughout ($+0.351$ gap). This confirms that learned policy --- not pre-authored schedules --- is responsible for emergent prosocial behavior.

    \item \textbf{Role masking contributes $+0.088$ to cooperation:} Disabling role-specific action masking (C6) reduces final cooperation from $0.851$ to $0.763$. The structured bias creates a cleaner learning signal for Q-learning agents.

    \item \textbf{System shocks reduce cooperation by $0.022$ but NPCs recover:} The full system (C1, $\cK = 0.851$) shows a dip-and-recover pattern around shock events compared to the no-shock condition (C5, $\cK = 0.873$). This demonstrates adaptive resilience --- shocks perturb but do not permanently degrade cooperation.

    \item \textbf{Dynamic $\lambda$ amplifies cooperation by $+0.158$:} C7 (static $\lambda = 0.325$) achieves $\cK = 0.693$, still showing emergence from the community reward signal, but the dynamic feedback loop (C1, $\cK = 0.851$) provides a $+0.158$ boost --- the second-largest single-component contribution after RL itself. This confirms the virtuous cycle mechanism: higher cooperation $\to$ higher $\lambda$ $\to$ more weight on community reward $\to$ even more cooperation.

    \item \textbf{LLM narration is more contextually grounded ($0.567$ vs.\ $0.488$) and uses $2.6\times$ richer vocabulary} than templates, while maintaining similar lexical coherence. Even a small local model (\texttt{qwen3:4b}) produces meaningfully better narrative output than the template fallback.
\end{itemize}

\subsection{What's Next}

With the initial automated experiments completed, the remaining future work focuses on scaling, validation, and real-world transfer:

\begin{itemize}
    \item \textbf{Scale the experiments:} Initial single-seed ablation is complete (Section~\ref{sec:preliminary_results}); next step is 100 bot playthroughs per condition with varied seeds for statistical significance
    \item Branching quest graphs with multiple concurrent quest lines to extend session duration beyond 200 turns
    \item Comparing tabular Q-learning vs.\ DQN for NPC behavior (trade interpretability for performance?)
    \item Player intent tracking via belief states (POMDP-style)
    \item Scaling to more NPCs with parameter sharing
    \item \textbf{User studies:} The bot experiments confirm the mechanisms work; human playtesting with post-session questionnaires is the next validation step
    \item Longer sessions ($>$500 turns) to investigate whether cooperation saturates or continues to evolve --- results show $\cK$ at $0.851$ after 500 turns, suggesting further room to grow
    \item More shock types and multi-shock interaction effects
\end{itemize}

\subsection{Real-World Transfer: Beyond Games}

The decomposed reward framework $R_i(t) = P_i(t) + G_i(t) + \lambda_i(t) \cdot C(t)$ is not game-specific. The individual--community tension it models appears naturally in at least two industrial domains where decentralised, self-organising control could replace costly centralised optimisation.

\subsubsection{Smart Factory Coordination}
In an automated production line, each cell (stamping, welding, painting, logistics) optimises its own KPI --- throughput, defect rate, idle time --- analogous to role-specific $G_i$ weights. The community signal $C(t)$ maps to plant-wide OEE (Overall Equipment Effectiveness). Supply-chain disruptions, equipment breakdowns, and demand spikes are direct analogues of famine, plague, and trade-boom shocks. Current Industry~4.0 approaches rely on centralised MPC controllers that must be re-modelled whenever the line changes. Our dynamic-$\lambda$ architecture offers a decentralised alternative: each cell runs a lightweight learner, the community reward propagates cooperation, and a new station joins the reward structure without replanning the entire plant. The ablation result C1~vs.~C3 ($+0.351$ cooperation from RL over fixed schedules) suggests learned policies would significantly outperform static dispatch rules.

\subsubsection{Automotive ECU Self-Tuning}
A modern engine's sensor--actuator subsystems (knock sensor, O$_2$ sensor, throttle body, MAP sensor, coolant temperature, crank position) map onto heterogeneous NPC roles: each maximises a local objective (detonation margin, fuel-trim accuracy, airflow response) while the community target $C(t)$ is engine-level efficiency (e.g.\ brake-specific fuel consumption). The dynamic $\lambda$ feedback loop enables \emph{continuous self-calibration}: a base map provides the pre-trained Q-table (cold-start), and on-road experience adjusts cooperation tendencies as altitude, fuel quality, and component wear shift the reward landscape. Component aging becomes a slow-onset shock; aftermarket modifications reset part of the state, and the system re-converges without manual tuning --- our C4 result ($\cK = 0.953$ with fast quest cycling) suggests rapid re-convergence is achievable.

\subsubsection{Bridging the Gap}
Two challenges must be addressed before transfer: (i)~\textbf{state-space scale} --- the 225-state tabular formulation must give way to function approximation (DQN or policy gradients) for high-dimensional, real-time sensor streams; and (ii)~\textbf{safety guarantees} --- the $P_i$ penalty term discourages but does not formally prevent catastrophic states. A hybrid architecture that enforces hard constraints from classical control theory while allowing RL-based optimisation within the safe envelope is the likely path forward.
